\subsection{The Module Graph}
\label{sec:module-import}

We begin with the coarsest level of analysis: module-level structures. Two graphs on the same vertex set $\mathcal{M}$ (the set of all \module{Mathlib} modules) capture complementary aspects of the library's organization. The \emph{module tree} $T$ (\S\ref{sec:module-names}--\S\ref{sec:namespace-tree}) inherits the cognitive classification that mathematicians have long applied to their subject. The \emph{module graph} $G_{\mathrm{module}}$ (\S\ref{sec:import-graph}) records the logical dependencies enforced by the compiler. Throughout this section, $T$ serves as the cognitive baseline against which we measure the structure of $G_{\mathrm{module}}$.

\paragraph{Modules}
\label{sec:module-names}
\begin{definition}[Module]\label{def:module}
A \emph{module} in \module{Mathlib} is identified by a dot-separated sequence of identifiers, such as $\module{Mathlib.Algebra.Group.Defs}$. Each module corresponds to a \texttt{.lean} source file, with dots replacing path separators. We write $\mathcal{M}$ for the set of all modules in \module{Mathlib}:
\begin{align*}
  \mathcal{M} = \{\; &\module{Mathlib.Init},\\
    &\module{Mathlib.Algebra.Group.Defs},\\
    &\module{Mathlib.Analysis.Normed.Group.Basic},\; \ldots\,\}.
\end{align*}
The \emph{depth} of a module $c_1.c_2.\cdots.c_k$ is~$k$, the number of dot-separated components.
\end{definition}

All statistics in this paper are computed from commit \texttt{534cf0b} (2 February 2026) of the \module{Mathlib} repository, using the official \texttt{importGraph} tool~\cite{importgraph}.\xinze{LeanScout pipeline: could cross-validate} At this commit, $|\mathcal{M}| = 7{,}563$. The set $\mathcal{M}$ serves as the shared vertex set for both graphs defined in this section. $\mathcal{M}$ contains only \module{Mathlib.*} modules; external dependencies (\texttt{Std}, \texttt{Batteries}, \texttt{Init}) are excluded, so the graphs capture \module{Mathlib}'s internal structure rather than the full dependency closure.

\begin{example}
\leavevmode
\begin{itemize}
  \item $\module{Mathlib}$ has depth~$1$.
  \item $\module{Mathlib.Algebra}$ has depth~$2$.
  \item $\module{Mathlib.Algebra.Group.Defs}$ has depth~$4$.
\end{itemize}
The maximum depth in \module{Mathlib} is~$7$.\footnote{Attained by, e.g., \module{Mathlib.CategoryTheory.Limits.Shapes.Pullback.IsPullback.Kernels}.}
\end{example}

\paragraph{The Module Hierarchy}
\label{sec:namespace-tree}

Modules are organized hierarchically: each dot-separated name mirrors a path in the source repository.

\begin{figure}[H]
\centering
\begin{subfigure}[t]{0.42\textwidth}
\centering
\begin{forest}
  for tree={
    font=\scriptsize\ttfamily,
    grow'=0,
    child anchor=west,
    parent anchor=south,
    anchor=west,
    calign=first,
    edge path={
      \noexpand\path [draw, \forestoption{edge}]
        (!u.south west) +(7.5pt,0) |- (.child anchor)\forestoption{edge label};
    },
    before typesetting nodes={
      if n=1 {insert before={[,phantom]}} {}
    },
    fit=band,
    before computing xy={l=13pt},
  }
[Mathlib/
  [Algebra/
    [Group/
      [Defs.lean]
    ]
  ]
  [Data/
    [Nat/
      [Notation.lean]
    ]
  ]
  [\ldots, edge={draw, densely dotted}]
]
\end{forest}
\caption{Directory structure}
\label{fig:dir-tree}
\end{subfigure}%
\hfill
\begin{subfigure}[t]{0.52\textwidth}
\centering
\begin{tikzpicture}[
  every node/.style={font=\scriptsize\ttfamily, inner sep=2pt, text=blue!60!black},
  dot/.style={circle, fill=blue!60!black, minimum size=3pt, inner sep=0pt},
  dep/.style={->, >=Stealth, thin, blue!60!black},
]
  \node[dot, label=above:{Mathlib}]              (root) at (1.8,0) {};
  \node[dot, label=left:{Algebra}]               (alg)  at (0,-1.3) {};
  \node[dot, label=right:{Data}]                 (dat)  at (3.6,-1.3) {};
  \node[dot, label=left:{Algebra.Group}]         (grp)  at (0,-2.6) {};
  \node[dot, label=right:{Data.Nat}]             (nat)  at (3.6,-2.6) {};
  \node[dot, label=below:{Algebra.Group.Defs}]   (gd)   at (0,-3.9) {};
  \node[dot, label=below:{Data.Nat.Notation}]    (nn)   at (3.6,-3.9) {};
  \draw[dep] (root) -- (alg);
  \draw[dep] (root) -- (dat);
  \draw[dep] (alg)  -- (grp);
  \draw[dep] (dat)  -- (nat);
  \draw[dep] (grp)  -- (gd);
  \draw[dep] (nat)  -- (nn);
  \node[font=\scriptsize, text=gray] at (1.8,-1.3) {$\cdots$};
\end{tikzpicture}
\caption{Module tree $T$: arrows point parent $\to$ child ($\in E_T$)}
\label{fig:namespace-tree}
\end{subfigure}
\caption{A fragment of the \module{Mathlib} source tree (commit \texttt{534cf0b}). Left: the directory layout. Right: the corresponding rooted tree~$T$. Each directory path maps to a module with dots replacing separators, e.g., \module{Mathlib/\allowbreak Algebra/\allowbreak Group/\allowbreak Defs.lean} $\leftrightarrow$ \module{Mathlib.\allowbreak Algebra.\allowbreak Group.\allowbreak Defs}.}
\label{fig:prefix-tree}
\end{figure}

We say $m$ is a \emph{prefix} of $m'$ if $m'$ extends $m$ by additional components. For example, \module{Mathlib.Algebra} is a prefix of \module{Mathlib.Algebra.Group.Defs}.

\begin{definition}[Module tree]\label{def:module-tree}
The prefix relation induces a rooted tree $T = (\mathcal{M}, E_T)$ with root $\module{Mathlib}$, which we call the \emph{module tree}, where
\[
  E_T = \{(m, m') \in \mathcal{M} \times \mathcal{M} \mid m \text{ is a prefix of } m' \text{ and } \mathrm{depth}(m') = \mathrm{depth}(m) + 1\}.
\]
Its leaves are modules (source files); its internal nodes are shared prefixes of modules, corresponding to directories in the source repository.
\end{definition} This hierarchy is distinct from Lean~4's \texttt{namespace} mechanism: in Lean~4, the \texttt{namespace} keyword creates a naming scope that need not align with file boundaries---a single namespace may span many files, and a single file may contribute declarations to multiple namespaces (\S\ref{sec:ns-mod-cross}). The two often coincide in practice, since \module{Mathlib}'s top-level directories (\module{Algebra}, \module{Analysis}, \module{CategoryTheory}, etc.) serve as both directory names and the primary organizational categories.

The top-level directories (children of the root) include \module{Algebra} (1257 files), \module{CategoryTheory} (982 files), \module{Analysis} (755 files), etc.

\paragraph{The Module Graph}
\label{sec:import-graph}

Each Lean source file begins with \texttt{import} statements that declare its dependencies. For example, the file \module{Mathlib/Data/Nat/Notation.lean} contains:

\begin{codebox}
\textsf{\bfseries\color{blue!60!black}Data/Nat/Notation.lean}\\[2pt]
\textbf{public import} \textcolor{blue!60!black}{Mathlib.Init}
\end{codebox}

\noindent
This file depends on exactly one other module: \module{Mathlib.Init}.

A file with more dependencies is \module{Mathlib/Algebra/Group/Defs.lean}, which imports nine modules (shown here in abbreviated form):

\begin{codebox}
\textsf{\bfseries\color{blue!60!black}Algebra/Group/Defs.lean}\\[2pt]
\textbf{public import} \textcolor{blue!60!black}{Batteries.Logic}\\
\textbf{public import} \textcolor{blue!60!black}{Mathlib.Algebra.Notation.Defs}\\
\textbf{public import} \textcolor{blue!60!black}{Mathlib.Algebra.Regular.Defs}\\
\hspace{1.5em}\textit{\color{gray}-- \ldots\ 6 more imports}
\end{codebox}

\noindent
Each file is a \emph{vertex} in the module graph; each \texttt{import} statement is a directed \emph{edge} from the importing module to the imported module.

Lean~4 supports both \texttt{import} and \texttt{public import}; the latter re-exports the imported module to downstream dependents. For the purposes of the module graph, both forms create an edge---they differ only in visibility, not in the dependency relationship.\xinze{CITE: Cite the Lean Reference Manual~\cite{lean4modules} here for the official definition of the module system and the \texttt{public import}/\texttt{import} distinction.}\oliver{MODULE: Oliver points out the module system (introduced Nov 2025) makes the \texttt{public import} vs \texttt{import} distinction structurally significant. Currently most Mathlib imports are \texttt{public}, but the community plans to gradually make imports private. This will fundamentally change transitive visibility and hence the structure of $G_{\mathrm{module}}$. We should: (1) note that our analysis treats both as equivalent edges, (2) acknowledge this is a simplification that will become increasingly invalid, (3) discuss what private imports would do to our redundancy and reachability metrics.}

\begin{definition}[Import set]
For each $m \in \mathcal{M}$, let $\mathrm{imports}(m) \subseteq \mathcal{M}$ denote the set of modules directly imported by $m$ (via either \texttt{import} or \texttt{public import}).
\end{definition}

\begin{definition}[Module graph]\label{def:module-graph}
The \emph{module graph} of \module{Mathlib} is $G_{\mathrm{module}} = (\mathcal{M}, E_{\mathrm{module}})$ where
\[
  E_{\mathrm{module}} = \{(m_1, m_2) \in \mathcal{M} \times \mathcal{M} \mid m_2 \in \mathrm{imports}(m_1)\}.
\]
\end{definition}

Since \module{Mathlib/Data/Nat/Notation.lean} contains \texttt{public import Mathlib.Init}, we have the edge $(\module{Data.Nat.Notation},\; \module{Init}) \in E_{\mathrm{module}}$:

\begin{figure}[H]
\centering
\begin{tikzpicture}[
  every node/.style={font=\small\ttfamily, inner sep=2pt, text=blue!60!black},
  dot/.style={circle, fill=blue!60!black, minimum size=3pt, inner sep=0pt},
  dep/.style={->, >=Stealth, thin, blue!60!black},
]
  \node[dot, label=above:{Init}]              (b) at (0,0) {};
  \node[dot, label=above:{Data.Nat.Notation}] (a) at (4,0) {};
  \draw[dep] (a) -- (b);
\end{tikzpicture}
\caption{The simplest import edge: \module{Data.Nat.Notation} imports \module{Init}. The arrow points from the importing module to the imported module. All names have the \module{Mathlib.}\ prefix removed.}
\label{fig:single-import-edge}
\end{figure}

\noindent\textbf{Note.}
The module tree $T$ and the module graph $G_{\mathrm{module}}$ are distinct structures on the same vertex set $\mathcal{M}$. In $T$, edges connect a module to its prefix parent in the directory hierarchy; in $G_{\mathrm{module}}$, edges reflect explicit \texttt{import} dependencies. Import edges need not respect the directory hierarchy.

The contrast between these two structures is conceptually significant. The module tree $T$ encodes a \emph{cognitive} organization of mathematics: its top-level directories---\module{Algebra}, \module{Analysis}, \module{Topology}---reflect a classification that mathematicians established long before formalization. The module graph $G_{\mathrm{module}}$, by contrast, encodes \emph{logical} dependencies that need not respect this classification. The extent to which $G_{\mathrm{module}}$ deviates from $T$---edges that cross directory boundaries, dependencies that link distant branches of the tree---reveals the internal structure of mathematics that no human taxonomy fully captures.

\begin{figure}[ht]
\centering
\begin{tikzpicture}[
  every node/.style={font=\small\ttfamily, inner sep=2pt, text=blue!60!black},
  dot/.style={circle, fill=blue!60!black, minimum size=3pt, inner sep=0pt},
  dep/.style={->, >=Stealth, thin, blue!60!black},
]
  \node[dot, label=above:{Init}]                (d) at (0,0) {};
  \node[dot, label=left:{Data.Int.Notation}]    (b) at (-1.5,-1.5) {};
  \node[dot, label=right:{Data.Nat.Notation}]   (c) at (1.5,-1.5) {};
  \node[dot, label=below:{Algebra.Group.Defs}]  (a) at (0,-3) {};

  \draw[dep] (a) -- (b);
  \draw[dep] (a) -- (c);
  \draw[dep] (b) -- (d);
  \draw[dep] (c) -- (d);
\end{tikzpicture}
\caption{A fragment of $G_{\mathrm{module}}$: \module{Algebra.Group.Defs} imports two modules that both depend on \module{Init}, forming a diamond. All names have the \module{Mathlib.}\ prefix removed.}
\label{fig:import-subgraph}
\end{figure}

The module graph $G_{\mathrm{module}}$ is DAG: we verified this by computing a topological ordering of all $7{,}563$ modules using the \texttt{importGraph} tool~\cite{importgraph}.\footnote{The \texttt{importGraph} tool outputs edges in the reverse direction (imported $\to$ importing). Our definition follows the standard dependency convention (importing $\to$ imported); acyclicity is invariant under edge reversal.} 

\paragraph{Transitive Reduction and Redundant Imports}
\label{sec:transitive-reduction}

Since $G_{\mathrm{module}}$ is a DAG, it has a unique \emph{transitive reduction}~\cite{ahu1972}: the smallest subgraph $G_{\mathrm{module}}^{-} \subseteq G_{\mathrm{module}}$ with the same reachability relation. An edge $(m_1, m_2) \in E_{\mathrm{module}}$ is \emph{transitively redundant} if $m_2$ is reachable from $m_1$ via a path of length $\ge 2$ in $G_{\mathrm{module}}$---that is, the dependency is already implied by other imports.

\begin{definition}[Transitive reduction]
The \emph{transitive reduction} of $G_{\mathrm{module}}$ is $G_{\mathrm{module}}^{-} = (\mathcal{M}, E_{\mathrm{module}}^{-})$ where
\[
  E_{\mathrm{module}}^{-} = E_{\mathrm{module}} \setminus \{(m_1, m_2)\in E_{\mathrm{module}} \mid \exists\, \text{path } m_1 \to \cdots \to m_2 \text{ in } G_{\mathrm{module}} \text{ of length} \ge 2\}.
\]
\end{definition}

\begin{definition}[Redundancy rate]
\label{def:redundancy-rate}
The \emph{redundancy rate} of $G_{\mathrm{module}}$ is the fraction of edges removed by transitive reduction:
\[
  r \;=\; \frac{|E_{\mathrm{module}} \setminus E_{\mathrm{module}}^{-}|}{|E_{\mathrm{module}}|}.
\]
\end{definition}

\begin{figure}[H]
\centering
\begin{subfigure}[t]{0.42\textwidth}
\centering
\vspace{0pt}
\begin{tikzpicture}[
  every node/.style={font=\scriptsize\ttfamily, inner sep=1pt, text=blue!60!black},
  dot/.style={circle, fill=blue!60!black, minimum size=3pt, inner sep=0pt},
  dep/.style={->, >=Stealth, thin, blue!60!black},
  redundant/.style={->, >=Stealth, orange!80!red, dashed, thick},
]
  \node[dot, label=above:{Init}]                (d) at (0,0) {};
  \node[dot, label=left:{Data.Int.Notation}]    (b) at (-1.2,-1.2) {};
  \node[dot, label=right:{Data.Nat.Notation}]   (c) at (1.2,-1.2) {};
  \node[dot, label=below:{Algebra.Group.Defs}]  (a) at (0,-2.4) {};

  \draw[dep] (a) -- (b);
  \draw[dep] (a) -- (c);
  \draw[dep] (b) -- (d);
  \draw[dep] (c) -- (d);
  \draw[redundant] (a) -- (d);
\end{tikzpicture}
\caption{Subgraph of $G_{\mathrm{module}}$}
\label{fig:tr-graph}
\end{subfigure}%
\hfill
\begin{subfigure}[t]{0.46\textwidth}
\vspace{0pt}\raggedleft
\scriptsize\ttfamily
\begin{tabular}[t]{@{}l@{}}
  \textnormal{\sffamily\bfseries\color{blue!60!black} Data/Int/Notation.lean} \\[1pt]
  \hspace{1em}public import \textcolor{blue!60!black}{Init} \\[6pt]
  \textnormal{\sffamily\bfseries\color{blue!60!black} Data/Nat/Notation.lean} \\[1pt]
  \hspace{1em}public import \textcolor{blue!60!black}{Init} \\[6pt]
  \textnormal{\sffamily\bfseries\color{blue!60!black} Algebra/Group/Defs.lean} \\[1pt]
  \hspace{1em}public import \textcolor{blue!60!black}{Data.Int.Notation} \\
  \hspace{1em}public import \textcolor{blue!60!black}{Data.Nat.Notation} \\
  \hspace{1em}\textcolor{orange!80!red}{public import \textcolor{blue!60!black}{Init}}
    \textnormal{\sffamily\color{orange!80!red} $\leftarrow$ would be redundant} \\
\end{tabular}
\caption{Source file headers}
\label{fig:tr-source}
\end{subfigure}
\caption{Illustrating transitive redundancy on the diamond from Figure~\ref{fig:import-subgraph}. If \module{Algebra.Group.Defs} were to import \module{Init} directly (\textcolor{orange!80!red}{dashed orange} edge), that edge would be transitively redundant: \module{Init} is already reachable via both \module{Data.Int.Notation} and \module{Data.Nat.Notation}. All names have the \module{Mathlib.}\ prefix removed.}
\label{fig:transitive-reduction}
\end{figure}

The raw module graph has $|E_{\mathrm{module}}| = 23{,}570$ edges. After transitive reduction, $|E_{\mathrm{module}}^{-}| = 19{,}448$ edges, removing $4{,}122$ redundant edges ($17.5\%$). These redundant imports are syntactically present in the source files but logically unnecessary: the Lean compiler resolves them through transitive dependencies regardless.\oliver{SHAKE: CRITICAL: Note here that Mathlib's CI runs the \texttt{shake} linter, which automatically removes certain types of redundant imports. The 17.5\% we measure is therefore the residual redundancy AFTER shake---imports that shake intentionally preserves (e.g., because they bring needed instances or notations into scope) or that shake does not target. This is a hard constraint on our data that must be stated explicitly in the methodology.} These redundant edges are not accidents of careless coding. Developers commonly import umbrella modules---such as importing all of group theory rather than a single lemma---to maintain a navigable mental model of the library. The redundancy reflects a cognitive need: the gap between what the compiler requires and what the human reader requires to understand a file's mathematical context.\oliver{SHAKE: Reframe this entire paragraph. Oliver objects to ``cognitive need'' and ``navigable mental model.'' The correct framing is \emph{development ergonomics}: developers keep these imports because (1) it makes future refactoring easier---if a file later needs a lemma from the umbrella, the import is already there; (2) it simplifies adding new theory to a file; (3) it ensures typeclass instances and notation remain available. This is about practical engineering convenience for library development, not about cognitive scaffolding.}

The \texttt{importGraph} tool applies transitive reduction by default when generating its output. We verified (via our own parser) that its reduced graph ($20{,}881$ edges including the umbrella root \module{Mathlib}) agrees exactly with our independently computed reduction on the shared $7{,}563$ modules: all $19{,}448$ edges in $G_{\mathrm{module}}^{-}$ appear in the tool's output, and the $1{,}433$ additional edges in the tool's output are precisely the edges from the umbrella module \module{Mathlib} to its direct children.

\paragraph{Module Containment Decay}
\label{sec:module-containment}

The module graph's relationship to the module hierarchy $T$ can be quantified at every depth of the hierarchy, not only at the top level. At each depth~$k$, we truncate both endpoints of every import edge to their first~$k$ dot-separated components and measure the proportion of edges whose truncated endpoints match.

\begin{definition}[Module containment at depth~$k$]
\label{def:module-containment}
For each depth $k \ge 1$, let $\mathrm{trunc}_k(m)$ denote the first~$k$ components of module~$m$. For example,
\[
  \mathrm{trunc}_2(\module{Mathlib.Algebra.Group.Defs}) = \module{Mathlib.Algebra}.
\]
The \emph{module containment ratio at depth~$k$} is
\[
  \mathrm{contain}_{\mathrm{mod}}(k) \;=\; \frac{|\{(m_1,m_2) \in E_{\mathrm{module}} \mid \mathrm{trunc}_k(m_1) = \mathrm{trunc}_k(m_2)\}|}{|E_{\mathrm{module}}|}.
\]
\end{definition}

Module depth is concentrated: $55.7\%$ of modules have depth~$4$ (e.g., \module{Mathlib.Data.Nat.Basic}), with a mean of~$4.22$ and a maximum of~$7$. Containment decays steeply: at depth~$1$ ($13$ top-level packages), $97.5\%$ of import edges stay within the same package; by depth~$4$, containment has fallen to~$8.8\%$ (Table~\ref{tab:module-containment} in Appendix~\ref{app:module-detail}). The steep collapse reveals that, within a discipline, dependencies routinely cross subdirectory boundaries. The containment curve for declaration-level dependencies (\S\ref{sec:containment-decay}) starts from a much lower baseline ($48.9\%$ at the top-level directory) but exhibits the same qualitative pattern.

\paragraph{Cross-Directory Dependencies}
\label{sec:cross-namespace}

\begin{definition}[Top-level directory]
For a module $m = \module{Mathlib.}X.Y.\cdots$, the \emph{top-level directory} is $\mathrm{dir}(m) = X$, the first directory component under the \module{Mathlib/} root.
\end{definition}

\begin{definition}[Intra-/inter-directory edge]
An edge $(m_1, m_2) \in E_{\mathrm{module}}$ is \emph{intra-directory} if $\mathrm{dir}(m_1) = \mathrm{dir}(m_2)$ and \emph{inter-directory} otherwise. For example, the edge
\[
  (\module{Mathlib.Algebra.Group.Defs},\;\; \module{Mathlib.Algebra.Notation.Defs})
\]
is intra-directory (both have $\mathrm{dir} = \module{Algebra}$), while
\[
  (\module{Mathlib.Algebra.Group.Defs},\;\; \module{Mathlib.Order.Defs})
\]
is inter-directory.
\end{definition}

\noindent
In $G_{\mathrm{module}}$, $14{,}464$ edges ($61.4\%$) are intra-directory and $9{,}106$ ($38.6\%$) are inter-directory. In $G_{\mathrm{module}}^{-}$, the proportions are similar: $12{,}226$ ($62.9\%$) intra and $7{,}222$ ($37.1\%$) inter.

The cross-directory heatmaps (Appendix~\ref{app:module-heatmaps}) confirm that while the diagonal dominates---most imports stay within the same top-level directory---significant off-diagonal traffic exists. The strongest inter-directory dependency is \module{RingTheory} $\to$ \module{Algebra} ($485$ edges), followed by \module{Data} $\to$ \module{Algebra} ($377$) and \module{Algebra} $\to$ \module{Data} ($318$).

Degree distribution, DAG depth, centrality, and robustness analyses for the module graph are reported in Appendix~\ref{app:module-detail}.

\noindent\textbf{Language contingency.}
\label{rem:language-contingency}
Note that the module graph's topology is shaped not only by mathematical dependencies but by Lean~4's specific compilation model: its prohibition of cyclic imports, its transitive visibility semantics, and the \texttt{public import} mechanism that allows selective re-export. Design choices such as fine-grained incremental compilation, selective re-export at the declaration level, or explicit interface files would alter the structure of $G_{\mathrm{module}}$ even for identical mathematical content.\xinze{CITE: Cite Huch and Wenzel~\cite{huch2024} here---they show how the Isabelle AFP's theory DAG structure enables distributed parallel build scheduling with $>$100x speedup, directly illustrating how compilation model shapes graph topology.} The structural statistics reported in this section---$17.5\%$ redundancy, DAG depth~$153$, mean out-degree~$3.1$---should be understood as properties of formalized mathematics \emph{in Lean~4}, not of formalized mathematics in general.\oliver{MODULE: This is the right place to discuss the module system explicitly. Oliver points out that the module system was introduced in Nov 2025 specifically to address the ``narrow waist'' / compilation bottleneck issue. The distinction between \texttt{public import} and \texttt{import} is a language-level mechanism designed to reshape exactly the graph we analyze. We should note: (1) our snapshot still has mostly \texttt{public} imports, (2) as private imports are adopted, transitive visibility will shrink, redundancy rates will change, and the funnel topology may flatten, (3) this makes our snapshot a baseline of the \emph{pre-module-system} era.}\oliver{SHAKE: Also note here that \texttt{shake} is another constraint shaping the 17.5\% figure.}

